\subsection{A fixed finite representation controls the fourth matching layer}

\begin{result}{E.130}{Proved}{Integral fourth-layer normalization to a fixed rational-point quotient}
Put \(N_i=\Pi^{(3)}/\Gamma_{i+1}\Pi^{(3)}\), let
\(t_0=5/16\), and fix its actual point section \(s_0\).
There is a characteristic finite quotient
\[
 Q=N_4/\overline{N_4^9(\Gamma_3N_4)^3}
 \quad\text{with}\quad
 1\longrightarrow M/3M\longrightarrow Q
       \longrightarrow H_9\longrightarrow1,
 \qquad M=\Gamma_3N_4,
 \tag{E.130.1}
\]
where \(M\simeq\mathbb Z_3^5\) and \(|Q|=3^{11}\).
Here \(N_4^9\) means the closed subgroup generated by ninth
powers, and \(H_9\) is the class-two Heisenberg quotient with
its three Hall coordinates reduced modulo nine.

For every pair in E.125, let \(s_2\) be its distinguished
support-matched pro-three class-two section, with the
literal reduction \(s_2\bmod9=s_0\bmod9\) supplied there.
There is a global integral class-four lift \(t\) of \(s_2\)
whose section in \(Q\) equals that of \(s_0\), after geometric
conjugacy. This allows the intervening class-three lift to
vary; it does not prescribe an arbitrary fixed class-three
lift.

More precisely, the two global \(M\)-valued cocycles lifting
a basis of the odd weight-three parameters can be chosen
so that their reductions modulo three are the reductions
of two fixed cocycles for the rational point \(s_0\).
The finite comparison is therefore uniform over the whole
family of E.125.
\end{result}
\begin{proof}
We first justify the finite quotient rather than treating
Hall-coordinate reduction as automatic. Choose an integral
Hall ordering with two generators in weight one, one in
weight two, two in weight three, and three in weight four.
The subgroup in (E.130.1) consists exactly of elements
whose first three coordinates are divisible by nine and
whose last five coordinates are divisible by three.

For this assertion, power-coordinate polynomials in the
power \(n\) are integer-valued of degree at most the
coordinate weight, and vanish at \(n=0\). Their expressions
in the basis \(\binom nj\) consequently have integral
coefficients. All \(\binom9j\), \(1\leq j\leq4\), are
divisible by three, and the first two are divisible by
nine. Thus every ninth power has the asserted coordinates.
The subgroup \(\Gamma_3N_4\) is abelian and its cubes give
exactly the last five coordinates divisible by three.
Collection preserves the asserted lattice. One direct
check uses
\[
 \binom{x+9}{k}-\binom{x}{k}
       =\sum_{j=1}^k\binom9j\binom{x}{k-j},
       \qquad k\leq4.
\]
It gives the required period nine in lower coordinates
for upper coordinates reduced modulo three; the lower
weight-two coordinate has period nine in its own
modulus. Upper coordinates occur linearly except for
products with lower weights, which also preserve
divisibility by three. Conversely the ordered powers of
the eight Hall generators, with these divisibilities,
belong to \(N_4^9(\Gamma_3N_4)^3\). This proves the
coordinate description, including the exact intersection
with \(M\), and hence (E.130.1).

Write \(c=[x,y]\). In rank two,
\(\Gamma_2N_4\) is abelian: modulo \(\Gamma_3\) it has
one generator, while
\([\Gamma_2,\Gamma_3]\subseteq\Gamma_5=1\).
It follows that \(c^3\) is central in \(Q\). It is also
fixed by the actual Galois action. Indeed its
weight-two coefficient changes by
\(3(\chi_3^2-1)\), divisible by nine, and all its upper
coordinates change by multiples of three.
Thus \(c^3\) is central even in the arithmetic extension
defined by \(Q\).

The module \(M\) has integral Lie coordinates in weights
three and four. To see integrality, write the free
logarithmic generators as
\(X=\log x\), \(Y=\log y\). In an expression of weight
at least three considered modulo weight five, only the
linear and quadratic parts of these logarithms occur.
The only denominator is two, a three-adic unit.
Set
\[
 U=[X,[X,Y]],\quad V=[Y,[X,Y]],\qquad
 P=[X,U],\quad Q_1=[Y,U]=[X,V],\quad R=[Y,V].
\]
These give an integral basis of \(M\). The actual
tangential action has
\(x\mapsto x^\chi,\ y\mapsto f^{-1}y^\chi f\) with
\(f\in\Gamma_2\). Its \(f\)-terms raise weight by two
and are invisible on \(M\). Twisting by a class-two
section with abelian cocycles \(a,b\) gives
\[
 g\cdot(u,v,z)=
 \left(\chi(g)^3u,\ \chi(g)^3v,\
 \chi(g)^4z+
 \chi(g)^3D(a(g),b(g))
                  \begin{pmatrix}u\\v\end{pmatrix}\right),
 \qquad
 D(a,b)=
 \begin{pmatrix}a&0\\ b&a\\0&b\end{pmatrix}.
 \tag{E.130.2}
\]
In particular the reduction of this module modulo three
is the same for every E.125 pair and for \(s_0\):
their abelian cocycles agree modulo nine.
Denote the rational-point module by \(M_0\).

There is a global integral class-four lift \(t'\) of
\(s_2\), by E.125 or E.66. Its class-three projection
is support matched, since the local weight-three first
cohomology vanishes at every support prime.
In \(Q\), \(t'\) and \(s_0\) have the same literal
projection to \(H_9\). Their difference is therefore
a continuous \(M_0/3M_0\)-valued one-cocycle \(\xi\).
Its projection to the two weight-three coordinates
is unramified outside three.

Here is a verification of this last assertion at every
type of place. At a fresh support prime, its exponent
\(n_j=(j+11)!\) is divisible by nine. The matched local
class-three point has inertia conjugate to \(x^{n_j}\)
or \(y^{n_j}\), which is the identity in the corresponding
quotient of \(Q\); \(s_0\) is unramified there.
At two, five, and eleven, the actual ordinary local
point has the same valuations as \(t_0\).
Since \(q-1\) is a three-adic unit, local units have
trivial pro-three Kummer class. A.31 identifies the full
local point classes. The class-three projection of
\(t'\) is consequently conjugate to \(s_0\) there.
Because their class-two representatives agree modulo
nine, the weight-one coordinates of the projected
conjugator are divisible by nine: divide their
Frobenius differences by \(q-1\).
The weight-two coordinate is divisible by three:
its remaining difference is divisible by nine and
\(v_3(q^2-1)=1\).
Its reduction in \(Q\) is thus a power of the arithmetic
central element \(c^3\), so it does not change the
finite section. The remaining weight-three
conjugator changes the odd cocycle by a coboundary;
its inertia value is zero because \(\chi_3\) is
unramified. Hence the odd projection of \(\xi\) has
zero inertia at these three places as well.

At every other \(q\ne3\), an integral section with
zero abelian inertia is unramified in every finite
nilpotent class. Indeed wild inertia maps trivially
to a pro-three group, and a first nonzero tame-inertia
coordinate in weight \(i\) would satisfy
\((q^i-q)u=0\) in a free \(\mathbb Z_3\)-module.
The nonzero integer \(q^i-q\) has zero kernel.
Induction on \(i\) proves the assertion. The actual
tangential action is unramified at \(q\), as retained
in A.31 and E.126. This completes the inertia check.

By E.67, the finite odd classes unramified outside
three are the two multiples of \(\kappa(3)\).
Choose once an integral generator
\(\gamma\in Z^1(G_{\mathbb Q},\mathbb Z_3(3))\)
whose reduction represents \(\kappa(3)\).
E.69 supplies two integral \(M\)-cocycles with
projections \(\gamma(1,0)\) and \(\gamma(0,1)\).
Subtract suitable multiples from \(t'\). The remaining
odd projection of \(\xi\) is a finite coboundary.
Lift its constant primitive to \(M\) and conjugate
accordingly. The remaining \(\xi\) is now a
\emph{literal} cocycle in the central
\(\mathbb F_3(4)^3\).

It is important to check ramification only after this
literal odd alignment. An \(M\)-coboundary can have
central inertia when \(M\) is ramified. We do not
assume those intermediate central terms vanish.
For the final pure central cocycle, tame invariance
at every support prime gives
\[
                 (q^{4-1}-1)\xi(T)=0
                          \quad\text{in }\mathbb F_3^3.
\]
Here \(q\equiv-1\pmod3\), so the coefficient is a unit
and \(\xi(T)=0\). At other \(q\ne3\), all original
sections and integral cocycles are unramified, and
the correcting \(M\)-coboundary is unramified because
the module action is. Thus this central cocycle is
unramified outside three. E.66 makes it the reduction
of an integral global weight-four cocycle.
Equivalently it is a combination of the three copies
of \(\eta_4=(\chi_3^4-1)/3\). Subtracting this combination
gives the asserted \(t\).

For the final uniformity assertion, choose two such
integral cocycles \(j_{0,1},j_{0,2}\) for \(M_0\), with
the same literal odd projections \(\gamma e_i\).
Choose integral lifts \(j_1,j_2\) for the actual module
and align those projections literally. Since the
module reductions agree, \(j_i-j_{0,i}\bmod3\) is a
genuine central finite cocycle. The preceding tame
argument makes it unramified at every support prime;
elsewhere outside three both modules and cocycles are
unramified. E.66 therefore permits an integral central
adjustment with
\[
                         j_i\bmod3=j_{0,i}\bmod3 .
 \tag{E.130.3}
\]
All adjustments are continuous cocycles or integral
geometric conjugacies. No finite cocycle with arbitrary
ramification is substituted for an integral lift.
\end{proof}

\begin{result}{E.131}{Proved}{One fixed finite Galois extension encodes every fourth-layer matching equation}
There is a fixed finite Galois extension
\(L/\mathbb Q\), unramified outside
\(\{2,3,5,11\}\), which encodes the complete
weight-four system of E.69 for every pair of E.125.
It can be defined by the actual rational-point
representation in \(Q\), the finite cusp-zero and
cusp-one path representations, the two reductions
\(j_{0,i}\bmod3\) from E.130, and \(\mu_9\).
At two, five, and eleven the constant vectors are zero
and the matrices are fixed. At a fresh support prime
the matrix and constant vector are functions of these
finite Frobenius data. They do not depend on later
prime choices or on the remaining ordinary local
unit coordinates.

Here are exact formulas in the integral Lie coordinates
of E.130. Write \(Z_i(g)\in\mathbb F_3^3\) for the
central component of \(j_{0,i}(g)\bmod3\), put
\[
 k_q=\frac{q^4-1}{3}\bmod3\in\mathbb F_3^\times ,
 \qquad Z(g)=\bigl(Z_1(g)\ Z_2(g)\bigr),
 \tag{E.131.1}
\]
and use the \(H_9\)-aligned Frobenius representative
already prescribed in E.125.
At a fresh \(A\)-support use the cusp-zero section;
at a fresh \(B\)-support use the cusp-one section.
The quotient of the rational-point Frobenius by that
cusp Frobenius belongs to \(M/3M\); denote its
weight-four component by \(P_{q,4}\).
Then the local data in E.69 can be chosen as
\[
 B_q=k_q^{-1}Z(F_q),\qquad
 d_q=-k_q^{-1}P_{q,4}
                     \quad(q\text{ fresh}),\qquad
 d_2=d_5=d_{11}=0 .
 \tag{E.131.2}
\]
The exact global matching condition is
\[
 (B_q-B_2)\lambda=d_q,\qquad
       \lambda\in\mathbb F_3^2,\quad
       q\in\Sigma\setminus\{2\}.
 \tag{E.131.3}
\]
This is an exact finite-representation description,
not an assertion that an explicit defining polynomial
for \(L\) or numerical values for all its Frobenius
classes have been computed.
\end{result}
\begin{proof}
Use the reference lift \(t\) normalized by E.130 and
the lifts \(j_i\) normalized by (E.130.3).
All finite representations defining \(L\) have finite
image and are continuous. The cusp data and actual
arithmetic action are unramified outside three by
E.128. The rational point and \(M_0\)-cocycles can
ramify only at two, three, five, and eleven.
For the latter assertion outside this set, inertia
acts trivially on \(M_0\), and the possible inertia
value of a cocycle would be fixed by \(q^{-1}F_q\).
Its eigenvalues on the two graded parts are
\(q^2\) and \(q^3\), neither equal to one in
\(\mathbb Q_3\); the resulting integral kernel is zero.
Thus \(L\) has the stated finite ramification.

For any support prime, the local odd cocycle has the
unique primitive
\[
                         u_q=\frac{\gamma(F_q)}{q^3-1}.
\]
Subtract its \(M\)-coboundary from \(j_i\). By
(E.130.2), its central Frobenius coordinate becomes
\[
 Z_i(F_q)-
       q^3D(a_0(F_q),b_0(F_q))\,u_q e_i
                           \pmod3 .
 \tag{E.131.4}
\]
This is unramified, and division by the unit \(k_q\)
identifies its local class with the global central
parameter. Formula (E.131.4) gives the fixed matrices
also at two, five, and eleven.
At every fresh support, E.125 prescribes
\(a_0(F_q)=b_0(F_q)=0\bmod9\), so the correction term
vanishes modulo three. This proves the formula for
\(B_q\).

We verify the constant term without relying on arbitrary
finite conjugacy. At an old support prime, the prescribed
point class is the full class of \(s_0\) by A.31.
The projected conjugator aligning it with \(s_2\)
has weight-one coordinates divisible by nine and
weight-two coordinate divisible by three, as in
E.130. Its image in \(Q\) is arithmetic central.
The remaining odd conjugator is divisible by three,
because \(t\) and \(s_0\) agree in \(Q\) and
\(q^3-1\) is a unit. It also acts trivially in \(Q\).
The even discrepancy consequently reduces to zero.
The map from integral local weight-four first
cohomology to finite first cohomology is injective:
its source has exponent three. Hence the actual
integral discrepancy is zero, proving the three
old constant terms.

At a fresh support, the local boundary description
represents the ordinary point by inertia \(x^{n_j}\)
or \(y^{n_j}\), transported to the chosen cusp path.
Its unit parameter can be removed integrally because
\(q-1\) is a three-adic unit. Its Frobenius is then
the chosen cusp Frobenius. Its inertia is the identity
in \(Q\). The class-two alignment again has coordinates
in nine and three respectively, and so acts trivially
in \(Q\).
The finite difference from \(t\) is exactly
\(P_q\in M/3M\).
Both cusp abelian Frobenius coordinates are zero.
Consequently (E.130.2) shows that eliminating the
odd component of \(P_q\) changes no weight-four
component: the off-diagonal term is zero.
The discrepancy of target from reference is therefore
\(-P_{q,4}\). Reduction detects the integral local
class, as above, and gives the second formula in
(E.131.2).

The formulas are independent of the auxiliary choices
of local lifts once the same arithmetic marking is
used on both sides. They describe local cohomology
classes; changing a Frobenius lift or an aligning
path changes representatives by the corresponding
coboundaries. In particular the chosen \(H_9\)-aligned
representatives can be used consistently in the finite
representation defining \(L\).
Pinning the unique central twist at two now gives
(E.131.3), exactly as in E.69.
The full local point class at a support depends only
on its two valuation coordinates, by A.31 and the
unit condition on \(q-1\). All fresh exponents are
divisible by nine, so the formulas use no later tail
or remaining local unit data.
\end{proof}

\begin{result}{E.132}{Proved}{A support-matched fourth-layer subfamily and an exact extension criterion for finite prefixes}
The family of E.125 has an infinite subfamily whose
distinguished class-two section lifts to a global
integral pro-three class-four section matched to the
ordinary local points at every support prime.
Every such section has a support-matched class-five
lift and a global integral class-six lift.

More generally, a legal finite E.125 prefix extends
to such a fourth-layer matched pair if and only if
the finite system (E.131.3) for its already specified
support, including two, five, and eleven, is soluble.
Every solubility witness can be retained while all
later primes and local roots are chosen. Thus an
inconsistent prefix cannot be repaired by a later
tail, and every consistent prefix has an infinite
continuation satisfying the entire fourth-layer
matching condition.

Matching here concerns the support primes, all
different from the coefficient prime three.
It does not assert local matching at three,
class-six matching, or a full pro-three section.
\end{result}
\begin{proof}
For the nonempty subfamily, at each new prime strengthen
E.125's finite \(H_9\) Frobenius prescription to the
following prescription in \(Q\): the actual point
section \(s_0\) has exactly the finite cusp-zero
Frobenius at an \(A\)-support, and exactly the finite
cusp-one Frobenius at a \(B\)-support. Include the
cusp-one path in the arithmetic data; its last
central coordinate must not be inferred from the
action on a preceding class.
E.128 proves that these finite marked conditions are
compatible with every old CRT condition and with
the required radical root. Only a finite three-power
precision may need enlargement. Infinitely many
new primes remain at every stage. The original
analytic and ordinary-limit argument of E.125 is
unchanged.

Now every \(P_{q,4}\) in (E.131.2) is zero.
The old constant terms are already zero.
The choice \(\lambda=0\), with central twist zero,
therefore solves the full system. E.130 and E.131
produce the actual integral matched class-four
section. E.66 supplies its next lift. Weight five
is odd and every \(q\in\Sigma\) is minus one modulo
three, so local weight-five first cohomology
vanishes. The lift is automatically support matched.
E.66 then supplies class six.

For the assertion about an arbitrary finite prefix,
let \(\lambda\in\mathbb F_3^2\) solve all its equations,
and choose any integral lifts of its two coordinates.
Put
\(\tau=-B_2\lambda\in\mathbb F_3^3\).
In the rational-point fibre choose the integral
class-four section
\[
       s_\lambda=s_0+
              j_{0,1}\lambda_1+j_{0,2}\lambda_2
                     +\eta_4\tau .
 \tag{E.132.1}
\]
The notation denotes the abelian \(M_0\)-torsor
operation, with the literal rational class-two
projection kept fixed.
Its finite geometric image over the needed
arithmetic data is the full group \(Q\).
Indeed inertia at five and eleven lies in the
kernel of that arithmetic data, which is unramified
outside three, and its two abelian valuations still
give a basis of \(Q/\Phi(Q)\).
The finite-group Frattini argument of E.126
therefore applies to this integral class-four
section just as it did to the point section.
E.128 explicitly permits this finite-class version.

At every later prime prescribe the finite cusp
Frobenius for \(s_\lambda\), rather than for \(s_0\).
Its \(H_9\) projection is still that of \(s_0\),
so these are extensions of the original legal
conditions. E.128 again supplies the radical-root
and CRT compatibility. Complete the E.125
construction with these choices.

For the resulting infinite pair, E.130 gives a
reference section reducing to \(s_0\) and integral
cocycles reducing to \(j_{0,i}\).
Apply the same two odd parameters and central
twist \(\tau\). The resulting global section
reduces to \(s_\lambda\) in \(Q\).
At each frozen prime it matches because the
assumed finite equations hold. At each later
prime it matches because the imposed cusp
Frobenius makes the corresponding weight-four
discrepancy zero. The actual integral conclusion
is exactly the finite-to-integral comparison
in E.131.

Conversely, the matrices and constant terms at
the frozen primes are unchanged by every later
tail. Any globally matched fourth-layer section
would solve those same finite equations by E.131.
If they are inconsistent, no continuation can
repair them. This proves both directions with
the finite prefix and its earlier local choices
held fixed.
\end{proof}
